\documentclass[11pt]{article}
\usepackage[a4paper,margin=25mm]{geometry}
\usepackage{amsmath,amssymb,amsthm,mathtools,booktabs,hyperref,xurl}
\setlength{\emergencystretch}{2em}
\newtheorem{theorem}{Theorem}
\newtheorem{lemma}[theorem]{Lemma}
\newtheorem{proposition}[theorem]{Proposition}
\newcommand{\Z}{\mathbb Z}\newcommand{\Q}{\mathbb Q}
\title{A half-trace square replaces the c-integrality predicate}
\author{\(\mathbb Z\mathbb Q\) Diophantine Project}
\date{24 November 2026}
\begin{document}
\maketitle
\begin{abstract}
An exact rational Pell classification turns one rational half-trace
square into a certificate of integrality for the source coordinate
\(c\). The required divisibility \(3\mid a\) already follows from
the integral coding prefix, before its final power recovery and without
assuming \(c\) integral. This permits the c-integer interface to be
replaced by the single half-trace row previously deleted in v101.
The argument treats both signs of \(a\) and the zero case.
Complete-table coordinate savings are separate compiler results;
this checkpoint asserts the arithmetic replacement theorem and its
source soundness and completeness.
\end{abstract}

\section{Only three possible negative-parameter seed orbits}

We use the proved rational classification of~\cite{Classification}.
For an integer \(a\), set
\[
 \Delta=a^2-1,\qquad D=4a-5,
 \qquad d^2-\Delta c^2=1,\qquad T=d-(a-2)c\in\Z.
\]
For \(|a|\ge2\), every solution is a signed translate of a
fundamental-interval seed under multiplication by
\(|a|+\sqrt\Delta\), and these translations preserve the exact
reduced denominator of \(c\). For \(a\ge3\), the nonintegral
orbit has denominator \(D/\gcd(D,3)\); at \(a=2\) all solutions
are integral.

The negative-parameter finite seed list admits the following stronger
description.

\begin{lemma}\label{negative}
Let \(a=-A\le-2\) and \(K=4A+5\). The complete set of distinct
fundamental-interval seeds is
\begin{equation}\label{two-seeds}
 (c_0,d_0)=(0,1),\qquad
 (c_2,d_2)=\left(\frac3K,\frac{5A+4}K\right),
\end{equation}
and, if and only if \(K=h^2\) for an integer \(h>0\), the
additional seed
\begin{equation}\label{third-seed}
                  (c_h,d_h)=\left(\frac2h,\frac{2A+1}h\right).
\end{equation}
\end{lemma}
\begin{proof}
The earlier classification says that seeds correspond precisely to
integers \(1\le t\le2A+1\) for which
\begin{equation}\label{seed-test}
                         v^2=(A^2-1)t^2+K
\end{equation}
is an integer square, with \(v>0\). The coordinates are
\[
 c_t=\frac{(A+2)t-v}K,\qquad
 d_t=\frac{(A+2)v-(A^2-1)t}K.
\]
If \(v>At\), then
\[
                  K-t^2=v^2-A^2t^2\ge2At+1.
\]
For \(t\ge3\), its left side is at most \(4A-4\), while
its right side is at least \(6A+1\), a contradiction. Thus
\(t=1\) or \(t=2\), giving respectively \(v=A+2\) and
\(v=2A+1\), and the two seeds in \eqref{two-seeds}.

If \(v\le At\), put \(j=At-v\). By \eqref{seed-test},
\(v>t\sqrt{A^2-1}\), so
\[
 0\le j<\frac{t}{A+\sqrt{A^2-1}}
       <\frac{2A+1}{2A-1}<2.
\]
Thus \(j\) is zero or one. At \(j=1\), equation
\eqref{seed-test} becomes
\[
                   (t+2)(t-2A-2)=0,
\]
contradicting the range of \(t\). At \(j=0\), it becomes
\(t^2=K\). Conversely, when \(K=h^2\), one has
\(1<h<2A+2\), so \(t=h\) is in range and gives
\eqref{third-seed}. Since \(K\ge13\), this positive integer
\(h\) is at least five and differs from one and two.
\end{proof}

\section{A square denominator forces integrality}

\begin{lemma}\label{square-denominator}
Suppose \(a\in3\Z\), \(c,d\in\Q\),
\(d^2-(a^2-1)c^2=1\), and \(T=d-(a-2)c\in\Z\).
If the reduced positive denominator of \(c\) is a square, then
\(c,d\in\Z\).
\end{lemma}
\begin{proof}
If \(a\ge3\), then \(3\nmid D\). A nonintegral solution
would have denominator \(D\), by the positive-parameter
classification. But \(D=4a-5\equiv3\pmod4\) is not a square.

If \(a=-A\le-3\), then \(3\mid A\) and
\(K=4A+5\equiv2\pmod3\) is nonsquare. Lemma~\ref{negative}
therefore leaves just the two seeds \eqref{two-seeds}. The nonintegral
seed has denominator \(K\), since \(3\nmid K\), and this
nonsquare denominator persists along its entire orbit.

At \(a=0\), the exact zero-case classification in~\cite{Classification}
gives
\[
                  c\in\{0,\pm1,\pm3/5,\pm4/5\}.
\]
The only nonintegral denominator is five, which is not a square.
These exhaust the multiples of three. In every case \(c\in\Z\),
and then \(d=(a-2)c+T\in\Z\).
\end{proof}

\begin{theorem}[One-square integrality certificate]\label{main}
Let \(a\in3\Z\), \(c,d,\zeta\in\Q\), and suppose
\begin{equation}\label{certificate}
 \begin{gathered}
 d^2-(a^2-1)c^2=1,\qquad T=d-(a-2)c\in\Z,\\
 \zeta^2=2\bigl(ad+(a^2-1)c+1\bigr).
 \end{gathered}
\end{equation}
Then \(c,d\in\Z\).
\end{theorem}
\begin{proof}
Write \(\kappa=a-2\). Substituting \(d=\kappa c+T\) into
the norm gives
\[
 Dc^2-2\kappa Tc-(T^2-1)=0,\qquad
 (Dc-\kappa T)^2=(a^2-1)T^2-D.
\]
The latter right side is an integer. Its rational square root is
therefore an integer, so \(Dc\in\Z\). This argument also
applies at \(a=0\). Let \(N\) be the reduced positive denominator
of \(c\); thus \(N\mid|D|\).

Set \(L_a=2a^2-2a-1\). The half-trace equation reads
\[
                      \zeta^2=2L_a c+2aT+2.
\]
The identity
\begin{equation}\label{gcd-identity}
                      8L_a+3=D(4a+1)
\end{equation}
shows that \(\gcd(L_a,D)\) divides three. Since \(3\mid a\),
the integer \(D\) is prime to three; it is also odd. Consequently
\(\gcd(2L_a,N)=1\), and the reduced denominator of
\(2L_a c+2aT+2\) is exactly \(N\). A rational square has a
square reduced denominator. Hence \(N\) is a square, and
Lemma~\ref{square-denominator} gives the conclusion.
\end{proof}

The hypothesis \(3\mid a\) cannot simply be omitted from this
signed theorem. The exact values
\[
 a=-5,\qquad c=-\frac3{25},\qquad d=-\frac{29}{25},
 \qquad T=-2,\qquad\zeta=\frac{14}5
\]
satisfy every equation in \eqref{certificate}, with nonintegral \(c\).
The retained source excludes this example by Lemma~\ref{prefix} below.

\section{The retained prefix supplies the arithmetic hypothesis}

\begin{lemma}\label{prefix}
For integers \(q,S,T_0,w,s\), define
\[
 n=q^{16},\qquad M=q^{32}=n^2,\qquad
 r=S(M-n)+T_0(M-1),\qquad a=(wM+1)rsM.
\]
Then \(3\mid Mr\), and hence \(3\mid a\).
\end{lemma}
\begin{proof}
Modulo three, \(n\) and \(M\) are both zero when \(3\mid q\)
and are both one otherwise. In the first case \(3\mid M\);
in the second \(3\mid M-n\) and \(3\mid M-1\), so
\(3\mid r\). No positivity or power-of-two hypothesis is used.
\end{proof}

The current source has exactly this integral prefix; its preliminary
reconstruction is given in~\cite[Section 2]{Signed}. The integer
inputs \(w,s,\gamma\) are retained. The integral recovery of
\(q,S,T_0\) uses the independent auxiliary power block and the
prefix equations, before the first-Pell index argument; it does not
use \(c\in\Z\). The first linear row is
\[
                   d=(a-2)c+bw+(4a-5)\gamma.
\]
It therefore gives precisely the integral shifted trace required by
Theorem~\ref{main}. Theorem~\ref{main} now recovers \(c,d\in\Z\)
before any remaining use of integral first-Pell classification.

In the specified interface, the retained explicit fixed-power circuit
makes \(q\) integral directly, and the same prefix lemma applies.

This order is noncircular. In particular no use of \(q=b^{5^{60}}\)
or of \(b\) being a power of two is needed to establish \(3\mid a\).
Both source signs are covered. The case \(a=0\), needed before
the source excludes \(s=0\), is included in the theorem.

\section{Canonical completeness and the quadratic interface}

The canonical positive source witnesses use \(a\ge2\),
\(B=2r+1\), and
\[
                     c=Y_B(a),\qquad d=X_B(a),
\]
where \(X_j+Y_j\sqrt{a^2-1}=(a+\sqrt{a^2-1})^j\).
The addition and doubling identities give
\[
 ad+(a^2-1)c=X_{B+1}(a)=X_{2r+2}(a)
                         =2X_{r+1}(a)^2-1.
\]
Hence
\begin{equation}\label{complete-root}
                            \zeta=2X_{r+1}(a)
\end{equation}
supplies the half-trace square. This is precisely the row removed in
v101; its canonical completion was already recorded
in~\cite[Section 6]{Signed}.

The replacement therefore preserves source-instance soundness and
completeness: remove the c-integer predicate, restore the half-trace
square, infer \(c\in\Z\) by Theorem~\ref{main}, and apply the
existing source argument. In the other direction, choose canonical
positive witnesses and \eqref{complete-root}, then apply simultaneous
integer-interface completeness to the reduced tuple of integer inputs.
The elliptic witnesses are chosen anew. All remaining interface and
coding premises remain in force.

With the already stored coordinate \(A_c=ac\), the restored row is
the quadratic polynomial
\[
                    \zeta^2-2ad-2aA_c+2c-2=0.
\]
A complete compiler must retain the definition of \(A_c\), and
replace any still-used register borrowed from the removed c chart.
In particular, if that chart supplied \(C_2=c^2\), retain one
ordinary quadratic square row for \(C_2\). Counting private
coordinates and common shared rows is a separate exact-table task.
No numerical bound follows from this manuscript alone.

\section{Reproducibility}

The independent checker verifies the symbolic denominator identity,
compares the three-seed formula with the original integer-square test
over a finite parameter range, checks square-denominator exclusions,
and checks exact prefix congruences and canonical half-trace completions.
The quantified statements are proved above. Run from the repository root:
{\small\begin{verbatim}
python -B \
  undecidability/pell_integrality_replacement_v106/verify_replacement.py
\end{verbatim}}

\begin{thebibliography}{9}
\small
\bibitem{Classification} Project, \emph{Rational Pell solutions with
an integral shifted trace}, reconstructed checkpoint v105,
24 November 2026.\par
\nolinkurl{undecidability/pell_classification_reconstructed_v105/pell_classification.tex}.
\bibitem{Signed} Project, \emph{Both first-Pell signs have the same
coding projection}, v101, 23 November 2026.\par
\nolinkurl{undecidability/pell_signed_branch_v101/signed_branch_projection.tex}.
\end{thebibliography}
\end{document}
